\documentclass[11pt,a4paper]{article}
\usepackage[T1]{fontenc}
\usepackage[utf8]{inputenc}
\usepackage{charter}
\usepackage[scaled=0.92]{helvet}
\usepackage[margin=24mm,top=22mm,bottom=20mm]{geometry}
\usepackage{xcolor}
\usepackage{booktabs}
\usepackage{titlesec}
\usepackage{microtype}

\definecolor{ink}{HTML}{15171C}
\definecolor{accent}{HTML}{7B1FA2}
\definecolor{muted}{HTML}{5B5F6B}
\color{ink}
\pagestyle{empty}
\setlength{\parskip}{5pt}
\setlength{\parindent}{0pt}

\begin{document}

\begin{center}
{\color{accent}\rule{\textwidth}{1.4pt}}\\[3mm]
{\LARGE\bfseries AEGIS: Adaptive Emergency Grid for Incident Steering\par}
\vspace{1.5mm}
{\normalsize An Intelligent Emergency Response \& Resource Optimization Platform\par}
\vspace{2mm}
{\color{accent}\rule{\textwidth}{1.4pt}}
\end{center}

\vspace{2mm}

Emergency dispatch is not primarily a routing problem; routing is a solved
commodity. The hard problems are allocation under scarcity, stability under
continuous change, and accountability for decisions affecting survival. AEGIS
is a sharded, event-driven backend built around those three.

Incidents are accepted durably on a region-partitioned event log, then triaged
by a versioned additive model that emits a per-term explanation rather than an
opaque rank. Allocation runs in two tiers over one shared cost function: a
sub-millisecond greedy admission decision, so no caller waits for a solver,
and a budgeted epoch that solves each regional shard \emph{exactly}. The key step,
requirement expansion, converts an NP-hard generalized assignment problem into
a rectangular linear assignment problem solvable optimally by
Jonker--Volgenant; scarcity, hospital capacity, coverage reserve
and plan stability enter as cost terms, refined by anytime Large Neighbourhood
Search. Re-optimisation is debounced and change-triggered, with switching costs
inside the objective so crews are never rerouted for marginal gain.

Conflict-freedom rests on three independent layers: single-writer sharding, TTL
reservation leases with all-or-nothing group acquisition, and optimistic
concurrency backed by a database constraint. Dependencies degrade through a
defined ladder rather than failing.

Evaluated on a replayed 12-hour national surge against a nearest-available
baseline: mean response $-14\%$, p90 $-39\%$, p99 $-25\%$, severity-weighted
$-18\%$, incidents never reached $-18\%$, utilisation $+37\%$, and zero
resource conflicts.

\vfill
{\color{muted}\footnotesize\hrule\vspace{2mm}
Full specification and reproducible evaluation harness accompany this submission.}

\end{document}
